\documentclass[xcolor=dvipsnames]{beamer} % dvipsnames gives more built-in colors
\usetheme{Singapore}
\useoutertheme{miniframes} % Alternatively: miniframes, infolines, split
\useinnertheme{circles}
\usecolortheme{dove}
\usepackage{hyperref}

\newcommand{\BE}{ \begin{enumerate}}
\newcommand{\EE}{ \end{enumerate}}
\newcommand{\BI}{ \begin{itemize}}
\newcommand{\EI}{ \end{itemize}}
\newcommand{\BEQ}{ \begin{equation}}
\newcommand{\EEQ}{ \end{equation}}
\newcommand{\BDM}{ \begin{displaymath}}
\newcommand{\EDM}{ \end{displaymath}}
\newcommand{\BFNS}{ \begin{footnotesize}}
\newcommand{\EFNS}{ \end{footnotesize}}
\newcommand{\BS}{ \begin{small}}
\newcommand{\ES}{ \end{small}}
\newcommand{\BSS}{ \begin{scriptsize}}
\newcommand{\ESS}{ \end{scriptsize}}
\newcommand{\BTi}{ \begin{tiny}}
\newcommand{\ETi}{ \end{tiny}}
\newcommand{\BVTi}{ \begin{verytiny}}
\newcommand{\EVTi}{ \end{verytiny}}
\newcommand{\balpha}{\mbox{\boldmath $\alpha$}}
\newcommand{\bOmega}{\mbox{\boldmath $\Omega$}}
\newcommand{\btau}{\mbox{\boldmath $\tau$}}
\newcommand{\bchi}{\mbox{\boldmath $\chi$}}
\newcommand{\bnabla}{\mbox{\boldmath $\nabla$}}
\newcommand{\grad}{\mbox{\boldmath $\nabla$}}
\newcommand{\boldeta}{\mbox{\boldmath $\eta$}}
\newcommand{\bphi}{\mbox{\boldmath $\varphi$}}
\newcommand{\ds}{\displaystyle}
\newcommand{\nl}{\ \\ }
\newcommand{\ud}{\textrm{ d}}
\newcommand{\bs}{\bigskip}
\newcommand{\BTc}{ \begin{tabular}{c}}
\newcommand{\BTl}{ \begin{tabular}{l}}
\newcommand{\BTr}{ \begin{tabular}{r}}
\newcommand{\ET}{ \end{tabular}}
\newcommand{\bu}{\mathbf{u}}
\newcommand{\bv}{\mathbf{v}}
\newcommand{\bx}{\mathbf{x}}
\newcommand{\be}{\mathbf{e}}
\newcommand{\bb}{\mathbf{b}}
\newcommand{\bk}{\mathbf{k}}
\newcommand{\bd}{\mathbf{d}}
\newcommand{\bn}{\mathbf{n}}
\newcommand{\bq}{\mathbf{q}}
\newcommand{\ba}{\mathbf{a}}
\newcommand{\bF}{\mathbf{F}}
\newcommand{\bI}{\mathbf{I}}
\newcommand{\bT}{\mathbf{T}}
\newcommand{\ust}{u_{*}}
\newcommand{\ddz}{\frac{\partial}{\partial z}}
\newcommand{\ddt}{\frac{\partial}{\partial t}}
\newcommand{\dudx}{\frac{\partial u}{\partial x}}
\newcommand{\dudz}{\frac{\partial u}{\partial z}}
\newcommand{\dvdz}{\frac{\partial v}{\partial z}}
\newcommand{\dudt}{\frac{\partial u}{\partial t}}
\newcommand{\dvdt}{\frac{\partial v}{\partial t}}
\newcommand{\ddudxx}{\frac{\partial^2 u}{\partial x^2}}
\newcommand{\dbudz}{\frac{\partial \bu}{\partial z}}
\newcommand{\dbudt}{\frac{\partial \bu}{\partial t}}
\newcommand{\partd}[2]{\frac{\partial #1}{\partial #2}}
\newcommand{\lagd}[2]{\frac{D #1}{D #2}}
\newcommand{\lagdh}[2]{\frac{D_h #1}{D #2}}

\newcommand{\arrow}[5]{ \put(#1,#2){\vector(#3,#4){#5}} }
\newcommand{\grid}{ 
\put(0,0){\vector(1,0){8}}
\put(0,1){\line(1,0){7}}  \put(0,2){\line(1,0){7}}
\put(0,3){\line(1,0){7}}  \put(0,4){\line(1,0){7}}
\put(0,5){\line(1,0){7}}
\put(0,0){\vector(0,1){6}}  
\put(1,0){\line(0,1){5}}  \put(2,0){\line(0,1){5}}  
\put(3,0){\line(0,1){5}}  \put(4,0){\line(0,1){5}}  
\put(5,0){\line(0,1){5}}  \put(6,0){\line(0,1){5}}  
\put(7,0){\line(0,1){5}}
\put(2.8,-0.7){$j$}
\put(7.3,-0.7){$x$}
\put(-0.7,2.8){$n$}
\put(-2.0,3.8){$n+1$}
\put(-0.7,5.3){$t$}}
\newcommand{\arakawa}{
\put(1,0){\line(0,1){4}}
\put(3,0){\line(0,1){4}}
\put(0,1){\line(1,0){4}}
\put(0,3){\line(1,0){4}}
}


\title{Numerical modelling of atmosphere and oceans \\ Lecture 5}
\author{P.L. Vidale \footnote{\tiny{based in part on J. Methven lecture notes from 2013}}}
\institute{
  National Centre for Atmospheric Science and Department of Meteorology\\
  University of Reading, UK\\
  p.l.vidale@reading.ac.uk
\begin{center}
	\begin{tabular}{cc}
     \includegraphics[width=0.15\textwidth]{Figures/AkioArakawa-203x250.jpg} &
     \includegraphics[width=0.15\textwidth]{Figures/Matsuno_photo.png}
	\end{tabular}
\end{center}
}
\date{Spring Term 2023}


\begin{document}
\frame{\titlepage}

\section[Outline]{}
\frame{\tableofcontents}

\section{Reminder of SWE}
\frame { \frametitle{Solving the SWEs numerically} 
	
	\BSS Need to chose finite difference
	methods to solve the {\bf shallow-water equations} on the
	plane tangent to the Earth's surface, {\em linearised} about a resting basic state:
	\begin{eqnarray}
	&&\frac{\partial \eta}{\partial t} + H \left(
	\frac{\partial u}{\partial x}+\frac{\partial v}{\partial y}\right) = 0,
	\\
	&&\frac{\partial u}{\partial t} - fv = - g\frac{\partial \eta}{\partial x}, \\
	&&\frac{\partial v}{\partial t} + fu = - g\frac{\partial \eta}{\partial y}, 
	\end{eqnarray}
	where $\eta$ is the surface elevation, $H$ is fluid depth, $(u,v)$ is
	the depth-averaged horizontal velocity, $f$ is the Coriolis parameter
	and $g$ is the gravitational acceleration. 
	
	
	
	Now discretise the problem using regular grids to describe $u$, $v$ and $\eta$ 
	
	e.g., $\eta_{ij}=$ free surface elevation at the $i$th longitude and $j$th
	latitude point.
	
	~
	
	Arakawa\footnote{\BTi Arakawa
		A. (1966) Computational design for long-term numerical integration of
		the equations of fluid motion: Two-dimensional incompressible
		flow. Part I. \emph{Journal of Computational Physics}, \textbf{1},
		119-143.\ETi} proposed a number of staggered finite difference
	grids that can be used to solve the shallow water equations.  
	
	\ESS }

\section{Arakawa's grids}
\frame
{
	\frametitle{Arakawa's Sm\"org{\aa}sbord}
	\BSS
	\begin{figure}[h]
		\centering
		\includegraphics[width=0.8\textwidth]{/Users/vidale/Documents/Presentations/"SWE 2D grid stagger types".pdf}
		\caption{Many of Arakawa's grid staggers}
		\label{fig:Arakawa_all}
	\end{figure}
	
	\ESS
	
}
\subsection{Real world examples}
\frame { \frametitle{Unified Model: New Dynamics vs EndGame grids} 
	
\BSS
The choice of grid stagger makes a difference even in very complex GCMs. Here is an example of "simply" switching the variables defined at the poles.
\begin{center}
	\begin{tabular}{cc}
		\includegraphics[width=0.75\textwidth]{Figures/New_Dynamics_vs_EndGame_grids.pdf}
	\end{tabular}
\end{center}

This is not at all a simple matter: implementation took years, but the consequences were profound in terms of scalability and numerical stability.

\ESS }


\frame { \frametitle{Unified Model DynCore evolution: stability and scalability} 
	
	\BSS
	The evolution from New Dynamics (2002) to EndGame (2014) meant greater scalability, so that we can use up to 12'000 cores efficiently, as well as numerical stability, which makes long climate simulations at high-resolution (up to N2560, 5km) feasible.
	
	\begin{center}
		\begin{tabular}{cc}
			\includegraphics[width=0.5\textwidth]{Figures/EndGame_stability.pdf}
			\includegraphics[width=0.5\textwidth]{Figures/EndGame_scalability_2.png}
		\end{tabular}
	\end{center}
		
	\ESS }


\section{Wave propagation in SWE}
\subsection{Propagation mechanisms}
\frame { \frametitle{Linear wave propagation} 

\BSS Why choose one grid staggering in preference to others? There are advantages and disadvantages to each configuration. Moreover, staggered grids represent wave propagation
differently.

~

The shallow water equations support {\em inertia-gravity waves} and  {\em Rossby waves}. 

~

{\bf Inertia-gravity waves}: propagation of anomalies in total fluid
depth. Wave restoration by gravitational acceleration. Typically much
higher frequency and smaller scale.

~

{\em Simplest form: seen as waves in free surface elevation on ocean with
	flat bottom. }

~

{\bf Rossby waves}: propagate by sideways displacement of potential
vorticity (PV) contours (where PV gradients exist). Wave restoration by
``Rossby elasticity''. Typically low frequency and large-scale.

~

{\em Simplest form: wave propagates westwards relative to a uniform
zonal flow, perpendicular to the \underline{poleward gradient} in planetary
vorticity.}

\ESS }

\subsection{A reminder of a few basic definitions}
\frame { \frametitle{Phase and Group velocities} 
	
	\BSS Wave movement can be described by two quantities:
		
	{\bf Phase velocity}: given in terms of the wavelength $\lambda$ and period $T$ as:
	\begin{equation}
	v_{p} ={\frac {\lambda }{T}} 
	\end{equation}
		
	Equivalently, in terms of the wave's angular frequency $\omega$, which specifies angular change per unit of time, and wavenumber (or angular wave number) $\kappa$, which represents the proportionality between the angular frequency $\omega$ and the linear speed (speed of propagation) $v_p$,
	
	\begin{equation}
	v_{p}={\frac {\omega }{k}}
	\end{equation}

	which means that a wave with frequency $\omega$ and wavenumber $\kappa$ moves to distance x after time t according to this relation: $\kappa x = \omega t$. The function $\omega(\kappa)$, which gives $\omega$ as a function of $\kappa$, is known as the \textbf{dispersion relation}. \\
	~
	~
	
	{\bf Group velocity}: of a wave is the velocity with which the overall shape of the waves' amplitudes, known as the modulation or envelope of the wave, propagates through space.
	
	\begin{equation}
	v_{g}\ \equiv \ {\frac {\partial \omega }{\partial k}}
	\end{equation}
	
	\ESS }



\section{Inertia-gravity waves}
\subsection{Analytic dispersion relation}

\frame
{
	\frametitle{Inertia-gravity waves}
	\BSS
One of the many types of wave motion supported by SWEs. Forcing comes from pressure (height) gradient and Coriolis. While the pure gravity waves are not dispersive ($\omega^2 = c_p^2 k^2$, where $c_p=\sqrt{gH}$), inertia-gravity waves are slighly dispersive, due to the $f_0$ term in the dispersion relation ($\frac{\omega^2} {f_0^2}=1 + R_D^2\left(k^2+l^2\right)$).
\begin{center}
	\begin{tabular}{cc}
		\includegraphics[width=0.45\textwidth]{Figures/Height_40.png}
	\end{tabular}
\end{center}

where $R_D = \sqrt{gH}/f_0$ is the \bf{Rossby deformation radius}.

	\ESS }

\frame
{
\frametitle{Analytical dispersion relation of inertia-gravity waves}
\BSS
Let us start by assuming a wave like solution of the form
\begin{equation}
\left \{ \begin{array}{c} \eta(x,y,t) \\ u(x,y,t) \\ v(x,y,t) \end{array} \right \} 
= 
\left \{ \begin{array}{c} \tilde{\eta} \\ \tilde{u} \\ \tilde{v} \end{array} \right \} e^{i (kx+ly-\omega t)}
\end{equation}
where $k$ and $l$ are the $x$ and $y$ components of the wave vector $\mathbf{k}$ and $\omega$ is the angular frequency of the wave and a constant Coriolis parameter ($f = f_0$). By plugging this solution into the model equations, we obtain the following system :
\[
\left( \begin{array}{ccc}
-i \omega & ikH & ilH \\
ikg & -i \omega & -f_0 \\
ilg & f_0 & -i \omega 
\end{array} \right )
\left( \begin{array}{c} \tilde{\eta} \\ \tilde{u} \\ \tilde{v} \end{array} \right )
= 
\mathbf{0}.
\]
For a non-trivial solution to exist, the determinant of this matrix
must equal zero, giving a {\em dispersion relation} between the wave
frequency and wave number: 
\BEQ \label{disp_exact}
\frac{\omega^2}{f_0^2} = 1 + R_D^2(k^2 + l^2), \EEQ 
where $R_D = \sqrt{gH}/f_0$ is the Rossby deformation radius. The
phase speed of the two inertia gravity waves is $c=\omega/\sqrt{k^2+l^2}
\approx \pm\sqrt{gH}$ and the Rossby wave has $\omega=0$. 

\ESS }

\frame
{
	\frametitle{A couple of notes on inertia-gravity waves and Rossby waves}
	\BSS
\begin{enumerate}
	\item On the previous page, oceanographers call the wave with $\omega=0$ a Rossby wave; we would simply call it geostrophic balance!
	\item Do try to solve the equations for the simpler case with $f_0=0$, as an exercise: what happens, and what do we call the waves?
	\item Can you remember the table in Lecture 4, with waves, names, properties? Start to fill it in. Again, an exercise...
	\item Watch these video to appreciate how fast some geophysical phenomena are: 1) \href{https://www.youtube.com/watch?v=oeKewmAoBEM}{1960 Chile Tsunami}; 2) \href{https://www.youtube.com/watch?v=jH3-hQjTGDQ}{2011 Japan Tsunami}. What waves are these?
	\item Before we go to the numerical solutions, please read the small "Simple note on advection" that is complementary to this lecture.
\end{enumerate}


\ESS }


\subsection{Inertia-gravity waves}
\frame
{
\frametitle{Dispersion relation for the A-grid on an $f$-plane}
\BSS
The same calculations can be performed on the space discretized equations obtained by using Arakawa's finite difference grids. When using the A-grid, we have the following spatially-discrete equations:
\begin{eqnarray*}
&&\frac{\ud \eta_{ij}}{\ud t} + H \left[ \frac{u_{i+1,j}-u_{i-1,j}}{2 d} + \frac{v_{i,j+1}-v_{i,j-1}}{2 d} \right] = 0,\\
&&\frac{\ud u_{ij}}{\ud t} - f_0 v_{ij} + g \frac{\eta_{i+1,j}-\eta_{i-1,j}}{2 d} = 0, \\
&&\frac{\ud v_{ij}}{\ud t} + f_0 u_{ij} + g \frac{\eta_{i,j+1}-\eta_{i,j-1}}{2 d} = 0, \\
\end{eqnarray*}
where we have assumed that $\Delta x = \Delta y = d$. By assuming again a wave-like solution of the form $u_{ij} = \tilde{u} e^{i(kx_i+ly_j-\omega t)}$, we obtain the following system:
\[
\left( \begin{array}{ccc}
-i \omega & iH\frac{\sin(kd)}{d} & iH\frac{\sin(ld)}{d} \\
ig\frac{\sin(kd)}{d} & -i \omega & -f_0 \\
ig\frac{\sin(ld)}{d} & f_0 & -i \omega 
\end{array} \right )
\left( \begin{array}{c} \tilde{\eta} \\ \tilde{u} \\ \tilde{v} \end{array} \right)
= 
\mathbf{0}
\]
and the corresponding dispersion relation is
\BEQ \label{disp_A}
\left(\frac{\omega^2}{f_0^2}\right)_A = 1 + \frac{R_D^2}{d^2} \left(\sin^2(kd) + \sin^2(ld)\right).
\EEQ
\ESS
}


\frame
{
\frametitle{Dispersion relations for the B- and C-grids}
\BSS
By doing the same for the B- and C-grids, we can obtain the following relations:
\begin{eqnarray}
\left(\frac{\omega^2}{f_0^2}\right)_B &=& 1 + 2\frac{R_D^2}{d^2} \left(1 - \cos(kd)\sin(ld)\right), \\
\left(\frac{\omega^2}{f_0^2}\right)_C &=& \cos^2(kd)\cos^2(ld) + 4\frac{R_D^2}{d^2} \left(\sin^2(\frac{kd}{2}) + \sin^2(\frac{ld}{2})\right).
\end{eqnarray}

These relations show that all three grids underestimate the wave
frequency. Also, for most of them, the slope of wave frequency with
respect to wavenumber (i.e., the group velocity) can be negative. This
means that energy might propagate in the wrong direction.  \ESS }


\frame
{
\frametitle{Graphical representation ($R_D/d = 10$), fine grid}
\begin{center}
\begin{tabular}{cc}
\includegraphics[width=0.45\textwidth]{Figures/gravity_exact_dispersion_Rd_10.eps} &
\includegraphics[width=0.45\textwidth]{Figures/gravity_Agrid_dispersion_Rd_10.eps}
\\
\includegraphics[width=0.45\textwidth]{Figures/gravity_Bgrid_dispersion_Rd_10.eps} &
\includegraphics[width=0.45\textwidth]{Figures/gravity_Cgrid_dispersion_Rd_10.eps}
\end{tabular}
\end{center}
\BSS Shading shows $\omega/f_0$ ranging from blue (low) to red (high)
(contour interval varies between panels). The graphs show solution for
$l=0$. On a fine grid ($d \ll R_D$, high-resolution) the C-grid gives the best results.  
\ESS
}


\frame
{
\frametitle{Graphical representation ($R_D/d = 1/10$), coarse grid}
\begin{center}
\begin{tabular}{cc}
\includegraphics[width=0.45\textwidth]{Figures/gravity_exact_dispersion_Rd_01.eps} &
\includegraphics[width=0.45\textwidth]{Figures/gravity_Agrid_dispersion_Rd_01.eps}
\\
\includegraphics[width=0.45\textwidth]{Figures/gravity_Bgrid_dispersion_Rd_01.eps} &
\includegraphics[width=0.45\textwidth]{Figures/gravity_Cgrid_dispersion_Rd_01.eps}
\end{tabular}
\end{center}
\BSS On a coarse (low-resolution) grid (i.e., when the grid size is larger than the
Rossby radius of deformation) the B-grid gives the best
results. Therefore, it was used in the old generation of global ocean
models.  \ESS }


\section{Excursus}
\subsection{Derivation of an important conservative quantity: PV}
\frame{
	\frametitle{Derivation of an important conservative quantity: PV}
	\BSS
	We saw in Lecture 4 how we can derive the Shallow Water Equations (SWEs). There are three independent variables: $u,v,p$ in that system:
	\begin{eqnarray}
	\partd{u}{x}+\partd{v}{y}+\partd{w}{z} & = & 0 \\
	\lagd{u}{t}-fv & = & -\frac{1}{\rho_o}\partd{p}{x} \\
	\lagd{v}{t}+fu & = & -\frac{1}{\rho_o}\partd{p}{y}
	\end{eqnarray}
	
	Using the vertically integrated hydrostatic balance: $-\frac{1}{\rho_o} \partd{p}{x} = -g \partd{h}{x}$, we can write: 
	\begin{eqnarray}
	\lagd{u}{t} & = & -g\partd{h}{x} +fv\\
	\lagd{v}{t} & = & -g\partd{h}{y} -fu
	\end{eqnarray}
	
	\ESS }

\frame{
	\frametitle{Derivation of an important conservative quantity: PV}
	\BSS
	It is possible to combine the three governing equations into a single one by cross-differentiation of the two momentum equations and substitution of mass continuity:
	\begin{eqnarray}
	&&\frac{\partial}{\partial y} \left( \lagd{u}{t} = -g\partd{h}{x} +fv \right) \\
	&&\frac{\partial}{\partial x} \left( \lagd{v}{t} = -g\partd{h}{y} -fu \right) 
	\end{eqnarray}
	
	Subtract the second from the first equation; use the definition of relative vorticity: $\zeta = \frac{\partial v}{\partial x} - \frac{\partial u}{\partial y}$:
	
	\begin{eqnarray}
	&&\frac{D_h}{Dt}\left({\zeta+f}\right) = -(\zeta+f) \left( \frac{\partial u}{\partial x} + \frac{\partial v}{\partial y} \right);\\
	&&\lagdh{h}{t} = -h \left( \frac{\partial u}{\partial x} + \frac{\partial v}{\partial y} \right)
	\end{eqnarray}
	to yield:
	\begin{eqnarray}
	\frac{D_h}{Dt}\left[\frac {\zeta+f}{h} \right]=0
	\label{PVcons}
	\end{eqnarray} 
	
	\ESS }

\subsection{Dispersion relation for Rossby waves on $\beta$-plane}
\frame{
	\frametitle{Dispersion relation for Rossby waves on $\beta$-plane}
	\BSS
	To derive the dispersion relation for Rossby waves we revert to the use of basic state and perturbations: $u=\bar{u}+u'; v=v'; \zeta = \bar{\zeta}+\zeta'$ on a $\beta$-plane: $f=f_0+\beta y$ in the barotropic vorticity equation: $\frac{D_h}{Dt}\left({\zeta+f} \right)=0$
	and making use of a streamfunction $\psi$ to define: $u'=-\partd{\psi'}{y}; v'=-\partd{\psi'}{x}$ to yield:
	\begin{eqnarray}
	\zeta'&=&\nabla^2\psi'\\
	\left(\partd{}{t}+\bar{u}\partd{}{x}\right)\nabla^2\psi'+\beta \partd{\psi'}{x}&=&0
	\end{eqnarray}
	
	and inserting a solution of this type: $\psi'=Re[\Psi exp(i\phi)]$, where $\phi=kx+ly-\omega t$ results in:
	
	\begin{eqnarray}
	(-\omega + k\bar{u})(-k^2-l^2)+k \beta & = & 0\\
	c_x-\bar{u} & = & -\frac{\beta}{K^2}
	\end{eqnarray}
	
	where: $\omega=\bar{u}k-\frac{\beta k}{K^2}$ and $K^2=k^2+l^2$.\\
	\ESS
}

\section{Rossby waves}
\subsection{Rossby waves}

\frame
{
	\frametitle{Rossby waves}
	\BSS
	One of the many types of wave motion supported by SWEs. Forcing comes from dependence of Coriolis term on latitude. They propagate \bf{westward}.
	\begin{center}
		\begin{tabular}{cc}
			\includegraphics[width=0.45\textwidth]{Figures/Height_480.png}
		\end{tabular}
	\end{center}
	
	\ESS }

\frame
{
\frametitle{Planetary (Rossby) waves (on a $\beta$-plane)}
\BSS
The restoring force of Rossby waves is the dependence of the Coriolis parameter on latitude. Rossby waves propagate information and energy westward across ocean basins. They are responsible for the \emph{westward intensification} associated with western boundary currents. The importance of Rossby waves for the large-scale circulation makes it important to choose an appropriate horizontal gridding scheme.

\vspace{0.2cm} To obtain the dispersion relation, we consider the linear shallow water equations expressed on a $\beta$-plane ($f=f_0 + \beta y$) with a constant depth. The dispersion relation, whose derivation will not be shown here\footnote{\BTi For details, see the book by A.E. Gill ``Atmosphere-Ocean Dynamics'', Academic press, 1982.\ETi}, is given by
\BEQ
\omega = - \frac{\beta k}{k^2 + l^2 + R_D^{-2}}.
\EEQ
\ESS
}


\frame
{
\frametitle{Dispersion relation for the A-, B- and C-grids}
\BSS
The discrete dispersion relations for Rossby waves on the first three Arakawa grids are:
\begin{eqnarray}
\left(\frac{\omega^2}{\beta d}\right)_A &=& \frac{-(R_D/d)^2 \sin(kd)\cos(ld)}{1 + (R_D/d)^2 \left[ \sin^2(kd) + \sin^2(ld) \right]},
\\
\left(\frac{\omega^2}{\beta d}\right)_B &=& \frac{-(R_D/d)^2 \sin(kd)}{1 + 2(R_D/d)^2 \left[1 - \cos(kd)\cos(ld) \right]},
\\
\left(\frac{\omega^2}{\beta d}\right)_C &=& \frac{-(R_D/d)^2 \sin(kd)\cos^2(ld/2)}{\cos^2(kd/2)\cos^2(ld/2) + 4(R_D/d)^2 \left[ \sin^2(kd/2) + \sin^2(ld/2) \right]}.
\end{eqnarray}
For these waves, the A-grid still gives the poorest results. The B- and C- grids are quite similar and give quite good results on a fine grid. When the grid is coarser, the errors for both increase. Unlike inertia-gravity waves, numerical methods can both under-estimate and over-estimate the analytical wave frequency.
\ESS
}


\frame
{
\frametitle{Graphical representation ($R_D/d = 10$), fine grid}
\begin{center}
\begin{tabular}{cc}
\includegraphics[width=0.45\textwidth]{Figures/rossby_exact_dispersion_Rd_10.eps} &
\includegraphics[width=0.45\textwidth]{Figures/rossby_Agrid_dispersion_Rd_10.eps}
\\
\includegraphics[width=0.45\textwidth]{Figures/rossby_Bgrid_dispersion_Rd_10.eps} &
\includegraphics[width=0.45\textwidth]{Figures/rossby_Cgrid_dispersion_Rd_10.eps}
\end{tabular}
\end{center}
\BSS
On a fine grid (high-resolution), the B- and C-grids give similar results.
\ESS
}


\frame
{
\frametitle{Graphical representation ($R_D/d = 1/10$), coarse grid}
\begin{center}
\begin{tabular}{cc}
\includegraphics[width=0.45\textwidth]{Figures/rossby_exact_dispersion_Rd_01.eps} &
\includegraphics[width=0.45\textwidth]{Figures/rossby_Agrid_dispersion_Rd_01.eps}
\\
\includegraphics[width=0.45\textwidth]{Figures/rossby_Bgrid_dispersion_Rd_01.eps} &
\includegraphics[width=0.45\textwidth]{Figures/rossby_Cgrid_dispersion_Rd_01.eps}
\end{tabular}
\end{center}
\BSS
Again, on a coarse grid (low-resolution), the B- and C-grids are quite similar for long waves.
\ESS
}


\section{Equatorial waves: propagation in SWEs}
\frame
{
	\frametitle{A general view of dispersion relations}
	\BSS
The review by Wheller and Nguyen (2015) goes through the derivations in the previous pages, and produces analytic solutions for a number of waves. The plot that shows the dispersion relation is a very useful reminder of various waves properties.

\begin{center}
\begin{tabular}{cc}
\includegraphics[width=0.45\textwidth]{Figures/Wheeler-Nguyen_Dispersion.png} &
\includegraphics[width=0.45\textwidth]{Figures/Wheeler-Nguyen_Dispersion-label.png}
\end{tabular}
\end{center}

\underline{Exercise}: extract information from this plot and insert in the Table provided at the start of Lecture 4.
\ESS
}


\section{Advanced topic: semi-implicit schemes}
\frame
{
\frametitle{Semi-implicit schemes}
\BSS



So far we have focussed on the spatial discretisation. Now let us
consider the discretisation of time. The two are linked, as you shall learn in Project 2.

The simplest choice for time derivatives is a simple forward or
centred (leapfrog) scheme. We can also choose the time-level of the RHS terms. For example,
\BEQ
\frac{u^{n+1}-u^n}{\Delta t}=fv^n ~~~~;~~~~~\frac{u^{n+1}-u^n}{\Delta t}=fv^{n+1}
\EEQ

are forward explicit and implicit schemes respectively.

~

Rule of thumb: implicit schemes are more stable than explicit
(e.g., implicit schemes for exponential decay are stable for any
time-step). However, fully implicit schemes for coupled equations are
difficult to solve and involve expensive iterations. 

~

Numerical weather prediction models use semi-implicit methods where
only the ``gravity wave terms'' are treated implicity and the rest are
explicit. The mass conservation equation can be {\em inverted} to
solve for future $\eta^{n+1}$. No iteration is required and schemes
are devised so that the matrix inversion is only done once. 

~

{\bf Benefit:} Lifts CFL restriction on time-step associated with
the fast GWs.

{\bf Cost:} Distorts and slows gravity waves.

\ESS
}

\frame
{
	\frametitle{The simplest possible approach to SI}
	\BSS
In atmospheric models, the fastest gravity waves, i.e., the external-gravity or “Lamb” waves, have speeds on the order of 300 $m s^{-1}$ , which is also the speed of sound. The typical time step for a model with a 10km mesh will thus have to be \underline{\hspace{1cm}}? This is unfortunate, because the external gravity modes are believed to play only a minor role in weather and climate dynamics. 

\begin{minipage}{0.4\textwidth}
	\begin{figure}[H]
		\includegraphics[trim={2cm 4.5cm 0cm 3cm},clip,width=4.cm]{/Users/vidale/Projects/Teaching/MTMW14/"2016 module"/Figures/SWE_GWs_Implicit}
		\label{fig:Spherical}
		\setbeamerfont{caption name}{size=\tiny}
		\caption{\label{fig:blue_rectangle} \tiny 1D SWEs in implicit form}
	\end{figure}
\end{minipage} \hfill
\begin{minipage}{0.55\textwidth}
Gravity waves are, therefore, commonly treated with implicit schemes, in order to mitigate this problem. However, this means solving a matrix problem (see following slides).	
\end{minipage}

\begin{minipage}{0.4\textwidth}
	\begin{figure}[H]
		\includegraphics[trim={1cm 2cm 0cm 2cm},clip,width=3.5cm]{/Users/vidale/Projects/Teaching/MTMW14/"2016 module"/Figures/SWE_GWs_FWBW_h_equation}\\
		\includegraphics[trim={0cm 2cm 0cm 2cm},clip,width=3.5cm]{/Users/vidale/Projects/Teaching/MTMW14/"2016 module"/Figures/SWE_GWs_FWBW_u_equation}
		\label{fig:Spherical}
		\setbeamerfont{caption name}{size=\tiny}
		\caption{\label{fig:blue_rectangle} \tiny 1D SWEs in FW-BW form}
	\end{figure}
\end{minipage} \hfill
\begin{minipage}{0.55\textwidth}
	Another approach is to go for the so called \emph{forward-backward} scheme, which eliminated the need for solving a matrix problem.
\end{minipage}

	\ESS
}


\frame
{
	\frametitle{A semi-implicit scheme for SWEs, part 1}
	\BSS

The shallow water equations linearised about a state of rest are below
discretised using a leapfrog scheme for Coriolis terms but a
trapezoidal scheme (mixed implicit-explicit) for the gravity wave
terms:
\begin{eqnarray*}
	\frac{u^{n+1}-u^{n-1}}{2\Delta t}-fv^n+\frac{g}{2}
	\left( \partd{h}{x}^{n+1} +\partd{h}{x}^{n-1} \right) & = & 0 \\
	\frac{v^{n+1}-v^{n-1}}{2\Delta t}+fu^n+\frac{g}{2}
	\left( \partd{h}{y}^{n+1} +\partd{h}{y}^{n-1} \right) & = & 0 \\
	\frac{h^{n+1}-h^{n-1}}{2\Delta t}+\frac{H}{2}
	\left( \partd{u}{x}^{n+1} +\partd{u}{x}^{n-1}
	+\partd{v}{y}^{n+1} +\partd{v}{y}^{n-1} 
	\right) & = & 0. 
\end{eqnarray*}

Re-arranging with future values on the left:
\begin{eqnarray*}
	u^{n+1}+\Delta t g \partd{h}{x}^{n+1} & = & A \\
	v^{n+1}+\Delta t g \partd{h}{y}^{n+1} & = & B \\
	h^{n+1}+\Delta t H \left( \partd{u}{x}^{n+1}+\partd{v}{y}^{n+1} \right) & = & C
\end{eqnarray*}


\ESS
}

\frame
{
	\frametitle{A semi-implicit scheme for SWEs, part 2}
	\BSS
Substituting $u^{n+1}$ and $v^{n+1}$ into the mass conservation
equation gives:
\begin{eqnarray*}
	\left\{ 1-\Delta t^2 gH \left( \partd{^2}{x^2}+\partd{^2}{y^2} \right) 
	\right\} h^{n+1}
	& = & C-\Delta t H \left\{ \partd{A}{x}+\partd{B}{y} \right\} \\
	{\cal L}\,h^{n+1} & = & F(u^{n-1},u^n,v^{n-1},v^n,h^{n-1},h^n) \\
	h^{n+1} & = & {\cal L}^{-1} F
\end{eqnarray*}

In words, the future depth can be found if the operator $\cal{L}$ can be
{\em inverted}. Once $h^{n+1}$ has been found, we can easily solve for
$u^{n+1}$ and $v^{n+1}$. 
	
	\ESS
}


\frame
{
	\frametitle{A semi-implicit scheme for SWEs, part 3}
	\BSS
The form of the operator $\cal{L}$ depends on the
representation of spatial derivatives by the numerical model. For
example, if a second order finite difference scheme is used:
\begin{equation}
	\partd{^2 h}{x^2}\approx \frac{h_{i+1}-2h_i+h_{i-1}}{\Delta x^2}
\end{equation}
then the $\cal{L}$ operator (in 1-D) becomes a tri-diagonal matrix:
\begin{equation}
	\left(
	\begin{array}{ccccccc}
		1-d & d & 0 & 0 & 0 & 0 & ... \\
		d & 1-2d & d & 0 & 0 & 0 & ... \\
		0 & d & 1-2d & d & 0 & 0 & ... \\
		& & \vdots & & & & ... 
	\end{array}
	\right) \left(
	\begin{array}{c}
		h^{n+1}_1 \\
		h^{n+1}_2 \\
		h^{n+1}_3 \\
		\vdots
	\end{array}
	\right) = \left(
	\begin{array}{c}
		F_1 \\
		F_2 \\
		F_3 \\
		\vdots
	\end{array}
	\right)
\end{equation}

where $d=gH\Delta t^2/\Delta x^2$. In this case the matrix is
time-invariant and only needs to be inverted once. The semi-implicit
scheme barely costs more that an explicit scheme per time-step but
enables a much longer time-step because it lifts the CFL restriction
associated with gravity wave speed $\sqrt{gH}$.
	
	\ESS
}

\frame
{
	\frametitle{QUESTIONS in preparation for P2 and test}
	\BSS
	\begin{enumerate}
		\item list one advantage and one disadvantage of using a A grid
		\item list one advantage and one disadvantage of using a C grid
		\item what happens with array dimensions when using a C grid?
		\item after dimensional analysis, you have opted for a C grid. How would you arrange the mass and velocity fields for simulating oceanic flow in a basin? Velocity at the boundaries? Mass at the boundaries? Why?
	\end{enumerate}
	\ESS
}

\frame
{
	\frametitle{Something to think about in preparation for Project 2}
	\BSS
	Waves in 2D can propagate in \underline{any} direction; they are not bound to travel along the zonal or meridional direction.
	
	\begin{center}
	\includegraphics[width=0.7\textwidth]{Figures/2DGaussian.png}
	\end{center}
	
	
	\underline{Question}: what does this 2D propagation mean for our CFL criterion? Is the time step requirement going to be stricter? Or is it going to be less strict?
	
	\medskip
	
	{\bf How to solve this puzzle:} draw a simple circular wave front on a 2D cartesian grid and think of the distance covered by the wave front over a sequence of 3-4 time steps. What happens to the signal along the $x$, $y$ axes and what happens along the diagonal?

	\ESS
}

\end{document}
